\chapter{1999年全国卷（理）}
\section{单选题}

\begin{problem}
如图，\(I\) 是全集，\(M\)，\(P\)，\(S\) 是 \(I\) 的 \(3\) 个子集，则阴影部分所表示的集合是（\quad）

\bitmapfigure[max width=0.38\linewidth]{img/1999/national_science/q1_fig1.png}

\choices
    {\((M\cap P)\cap S\)}
    {\((M\cap P)\cup S\)}
    {\((M\cap P)\cap \overline{S}\)}
    {\((M\cap P)\cup \overline{S}\)}
\end{problem}

\begin{answer}
C
\end{answer}

\begin{solution}
图中阴影内任意元素都属于 \(M\) 和 \(P\)，因此先取交集 \(M\cap P\)；阴影不在 \(S\) 内，因此还要取 \(S\) 的补集. 所以阴影部分表示
\((M\cap P)\cap\overline{S}\)，故选 C.
\end{solution}

\begin{problem}
已知映射 \(f:A\to B\)，其中，集合 \(A=\{-3,-2,-1,1,2,3,4\}\)，集合 \(B\) 中的元素都是 \(A\) 中元素在映射 \(f\) 下的象，且对任意的 \(a\in A\)，在 \(B\) 中和它对应的元素是 \(|a|\)，则集合 \(B\) 中元素的个数是（\quad）

\choices
    {\(4\)}
    {\(5\)}
    {\(6\)}
    {\(7\)}
\end{problem}

\begin{answer}
A
\end{answer}

\begin{solution}
由 \(f(a)=|a|\)，逐一取集合 \(A\) 中元素的象，得到
\(f(-3)=3\)，\(f(-2)=2\)，\(f(-1)=1\)，\(f(1)=1\)，\(f(2)=2\)，\(f(3)=3\)，\(f(4)=4\).
因此
\(B=\{f(a):a\in A\}=\{1,2,3,4\}\)，共有 \(4\) 个不同元素，故选 A.
\end{solution}

\begin{problem}
若函数 \(y=f(x)\) 的反函数是 \(y=g(x)\)，\(f(a)=b\)，\(ab\neq 0\)，则 \(g(b)\) 等于（\quad）

\choices
    {\(a\)}
    {\(a^{-1}\)}
    {\(b\)}
    {\(b^{-1}\)}
\end{problem}

\begin{answer}
A
\end{answer}

\begin{solution}
反函数的定义给出 \(g(f(x))=x\)（在 \(f\) 的定义域内）. 由题设 \(f(a)=b\)，代入得
\(g(b)=g(f(a))=a\)，故选 A. 条件 \(ab\neq0\) 不影响这一反函数关系.
\end{solution}

\begin{problem}
函数 \(f(x)=M\sin(\omega x+\varphi)\)（\(\omega>0\)）在区间 \([a,b]\) 上是增函数，且 \(f(a)=-M\)，\(f(b)=M\)，则函数 \(g(x)=M\cos(\omega x+\varphi)\) 在 \([a,b]\) 上（\quad）

\choices
    {是增函数}
    {是减函数}
    {可以取得最大值 \(M\)}
    {可以取得最小值 \(-M\)}
\end{problem}

\begin{answer}
C
\end{answer}

\begin{solution}
由 \(f(a)=-M<f(b)=M\) 且 \(f\) 在 \([a,b]\) 上为增函数可知 \(M>0\). 令 \(u(x)=\omega x+\varphi\)，并记 \(u_a=u(a)\)、\(u_b=u(b)\)，则 \(u\) 随 \(x\) 增大而严格增大，且
\(\sin u_a=-1\)，\(\sin u_b=1\).

因为 \(M\omega>0\)，且 \(f\) 在整个区间上递增，所以在相位变化的区间内必须有
\(f'(x)=M\omega\cos u(x)\geqslant0\)，即 \(\cos u(x)\geqslant0\). 又由 \(\sin u(a)=-1\)、\(\sin u(b)=1\)，并结合 \(u(x)\) 严格递增，只能取
\(u(a)=-\frac{\pi}{2}+2k\pi\)，\(u(b)=\frac{\pi}{2}+2k\pi\)，其中 \(k\in\Z\). 因而区间 \([u(a),u(b)]\) 包含 \(2k\pi\). 取 \(x_0\in[a,b]\) 使 \(u(x_0)=2k\pi\)，则
\(g(x_0)=M\cos(2k\pi)=M\). 又因 \(|\cos u|\leqslant1\) 且 \(M>0\)，\(g(x)\leqslant M\) 对所有 \(x\in[a,b]\) 成立，所以 \(g\) 可以取得最大值 \(M\)，故选 C.
\end{solution}

\begin{problem}
若 \(f(x)\sin x\) 是周期为 \(\pi\) 的奇函数，则 \(f(x)\) 可以是（\quad）

\choices
    {\(\sin x\)}
    {\(\cos x \)}
    {\(\sin 2x \)}
    {\(\cos 2x \)}
\end{problem}

\begin{answer}
B
\end{answer}

\begin{solution}
设 \(h(x)=f(x)\sin x\)，逐项检验：
当 \(f(x)=\sin x\) 时，\(h(x)=\sin^2x\)，满足 \(h(-x)=h(x)\)，是偶函数；
当 \(f(x)=\cos x\) 时，
\(h(x)=\sin x\cos x=\frac{1}{2}\sin 2x\)，满足 \(h(-x)=-h(x)\)，且 \(h(x+\pi)=h(x)\)，所以它是周期为 \(\pi\) 的奇函数；
当 \(f(x)=\sin2x\) 时，
\(h(x)=\sin2x\sin x=\frac{1}{2}(\cos x-\cos3x)\)，是偶函数；
当 \(f(x)=\cos2x\) 时，
\(h(x)=\cos2x\sin x=\frac{1}{2}(\sin3x-\sin x)\) 虽是奇函数，但
\(h(\frac{\pi}{6})=\frac{1}{4}\)，\(h(\frac{7\pi}{6})=-\frac{1}{4}\)，所以 \(h(x+\pi)\) 不恒等于 \(h(x)\)，不以 \(\pi\) 为周期. 因此只有选项 B 符合.
\end{solution}

\begin{problem}
在极坐标系中，曲线 \(\rho=4\sin\left(\theta-\frac{\pi}{3}\right)\) 关于（\quad）

\choices
    {直线 \(\theta = \frac{\pi}{3}\) 轴对称}
    {直线 \(\theta = \frac{5\pi}{6}\) 轴对称}
    {点 \(\left(2, \frac{\pi}{3}\right)\) 中心对称}
    {极点中心对称}
\end{problem}

\begin{answer}
B
\end{answer}

\begin{solution}
将极坐标方程乘以 \(\rho\)，并利用 \(x=\rho\cos\theta\)、\(y=\rho\sin\theta\)，得
\(\rho^2=4\rho\sin\left(\theta-\frac{\pi}{3}\right)=2y-2\sqrt{3}x\)，所以
\(x^2+y^2=2y-2\sqrt{3}x\)，即
\((x+\sqrt{3})^2+(y-1)^2=4\).
该圆的圆心为 \((-\sqrt{3},1)\)，其极角为 \(\frac{5\pi}{6}\)，所以圆关于直线 \(\theta=\frac{5\pi}{6}\) 轴对称，故选 B.
\end{solution}

\begin{problem}
将若干毫升水倒入底面半径为 \(2 \mathrm{~cm}\) 的圆柱形器皿中，量得水面的高度为 \(6 \mathrm{~cm}\)，若将这些水倒入轴截面是正三角形的倒圆锥形器皿中，则水面的高度是（\quad）

\choices
    {\(6\sqrt{3} \mathrm{~cm}\)}
    {\(6 \mathrm{~cm}\)}
    {\(2\sqrt[3]{18} \mathrm{~cm}\)}
    {\(3\sqrt[3]{12} \mathrm{~cm}\)}
\end{problem}

\begin{answer}
B
\end{answer}

\begin{solution}
圆柱形器皿中水的体积为 \(V=\pi\cdot 2^2\cdot 6=24\pi\ \mathrm{cm}^3\). 设倒圆锥中水面的高度为 \(h\)，水面半径为 \(r\). 由于倒圆锥的轴截面是正三角形，轴截面的底边为 \(2r\)，高为 \(h\)，所以 \(h=\sqrt{3}r\)，即 \(r=h/\sqrt{3}\). 因此水的体积为 \(\frac{1}{3}\pi r^2h=\frac{\pi h^3}{9}\). 由体积不变得 \(\frac{\pi h^3}{9}=24\pi\)，所以 \(h^3=216\). 因为 \(h>0\)，故 \(h=6\ \mathrm{cm}\)，所以选 B.
\end{solution}

\begin{problem}
若 \(\left(2x + \sqrt{3}\right)^4 = a_0 + a_1x + a_2x^2 + a_3x^3 + a_4x^4\)，则 \((a_0 + a_2 + a_4)^2 - (a_1 + a_3)^2\) 的值为（\quad）

\choices
    {\(1\)}
    {\(-1\)}
    {\(0\)}
    {\(2\)}
\end{problem}

\begin{answer}
A
\end{answer}

\begin{solution}
设 \(P(x)=(2x+\sqrt{3})^4\)，并记偶次项系数和为 \(E=a_0+a_2+a_4\)，奇次项系数和为 \(O=a_1+a_3\). 则
\(E=\frac{P(1)+P(-1)}{2}\)，\(O=\frac{P(1)-P(-1)}{2}\)，从而
\[
E^2-O^2=(E+O)(E-O)=P(1)P(-1)=(2+\sqrt{3})^4(\sqrt{3}-2)^4=1
\]
故选 A.
\end{solution}

\begin{problem}
直线 \(\sqrt{3}x+y-2\sqrt{3}=0\) 截圆 \(x^{2}+y^{2}=4\) 得的劣弧所对的圆心角为（\quad）

\choices
    {\(\frac{\pi}{6}\)}
    {\(\frac{\pi}{4}\)}
    {\(\frac{\pi}{3}\)}
    {\(\frac{\pi}{2}\)}
\end{problem}

\begin{answer}
C
\end{answer}

\begin{solution}
圆心 \(O(0,0)\) 到直线 \(\sqrt{3}x+y-2\sqrt{3}=0\) 的距离为
\(d=\frac{2\sqrt{3}}{\sqrt{3+1}}=\sqrt{3}\).
设劣弧所对圆心角的一半为 \(\theta\)，则在由圆心、弦中点和弦端点组成的直角三角形中，
\(\cos\theta=\frac{d}{2}=\frac{\sqrt{3}}{2}\)，故 \(\theta=\frac{\pi}{6}\)，劣弧所对的圆心角为 \(2\theta=\frac{\pi}{3}\)，故选 C.
\end{solution}

\begin{problem}
如图，在多面体 \(ABCDEF\) 中，已知面 \(ABCD\) 是边长为 \(3\) 的正方形，\(EF\parallel AB\)，\(EF=\frac{3}{2}\)，\(EF\) 与面 \(AC\) 的距离为 \(2\)，则该多面体的体积为（\quad）

\bitmapfigure[max width=0.42\linewidth]{img/1999/national_science/q10_fig1.png}

\choices
    {\(\frac{9}{2}\)}
    {\(5\)}
    {\(6\)}
    {\(\frac{15}{2}\)}
\end{problem}

\begin{answer}
D
\end{answer}

\begin{solution}
由于 \(EF\parallel AB\)，且 \(AB\) 与 \(AC\) 相交，\(EF\) 与 \(AC\) 的公垂线垂直于底面 \(ABCD\)，所以题设距离 \(2\) 就是该多面体的高. 设从底面起高为 \(2t\)（\(0\leqslant t\leqslant1\)）的水平截面面积为 \(S_t\). 该截面是底面正方形按比例 \(1-t\) 缩小后与线段 \(EF\) 按比例 \(t\) 叠加得到的矩形，且其两条边分别平行于 \(EF\) 及垂直于 \(EF\). 因此，平行于 \(EF\) 的边长由底面的 \(AB=3\) 线性变到顶边的 \(EF=\frac{3}{2}\)，为

\[
3(1-t)+\frac{3}{2}t
\]

另一条边长由底面的 \(AD=3\) 线性变到顶端的 \(0\)，为 \(3(1-t)\). 所以

\[
S_t=\left(3(1-t)+\frac{3}{2}t\right)3(1-t)
\]

令实际高度为 \(z=2t\)，则 \(\mathrm dz=2\,\mathrm dt\)，由截面积积分得
\[
V=\int_0^2 S_{z/2}\,\mathrm dz
 =2\int_0^1\left(3-\frac{3}{2}t\right)3(1-t)\,\mathrm dt
 =2\int_0^1\left(9-\frac{27}{2}t+\frac{9}{2}t^2\right)\,\mathrm dt
 =2\left(9-\frac{27}{4}+\frac{3}{2}\right)
 =\frac{15}{2}
\]
所以该多面体的体积为 \(\frac{15}{2}\)，故选 D.
\end{solution}

\begin{problem}
若 \(\sin \alpha>\tan \alpha>\cot \alpha\)（\(-\frac{\pi}{2}<\alpha<\frac{\pi}{2}\)），则 \(\alpha\in\)（\quad）

\choices
    {\(\left(-\frac{\pi}{2}, -\frac{\pi}{4}\right)\)}
    {\(\left(-\frac{\pi}{4}, 0\right)\)}
    {\(\left(0, \frac{\pi}{4}\right)\)}
    {\(\left(\frac{\pi}{4}, \frac{\pi}{2}\right)\)}
\end{problem}

\begin{answer}
B
\end{answer}

\begin{solution}
在给定区间内 \(\cos\alpha>0\)，且 \(\tan\alpha\) 有意义. 令 \(t=\tan\alpha\)，则
\(\sin\alpha=\frac{t}{\sqrt{1+t^2}}\).
由
\(\sin\alpha-\tan\alpha=t\left(\frac{1}{\sqrt{1+t^2}}-1\right)>0\)
可知 \(t<0\). 在 \(t<0\) 时，\(\tan\alpha>\cot\alpha\) 等价于 \(t>1/t\)，两边同乘负数 \(t\) 后不等号反向，得 \(t^2<1\). 因而
\(-1<t<0\)，从而 \(-\frac{\pi}{4}<\alpha<0\)，故选 B.
\end{solution}

\begin{problem}
如果圆台的上底面半径为 \(5\)，下底面半径为 \(R\)，中截面把圆台分为上，下两个圆台，它们的侧面积的比例为 \(1:2\)，那么 \(R=\)（\quad）

\choices
    {\(10\)}
    {\(15\)}
    {\(20\)}
    {\(25\)}
\end{problem}

\begin{answer}
D
\end{answer}

\begin{solution}
设圆台母线长为 \(l\)，中截面的半径为 \(r=\frac{R+5}{2}\). 中截面平行于两底面且位于母线中点，所以分成的上、下两个圆台母线长均为 \(l/2\). 由圆台侧面积公式 \(S=\pi(r_1+r_2)l\)，侧面积之比等于相应上下底面半径和之比：
\[
\frac{5+r}{r+R}=\frac{1}{2}
\]
代入 \(r=\frac{R+5}{2}\)，得
\[
\frac{R+15}{3R+5}=\frac{1}{2}
\]
所以 \(R=25\)，故选 D.
\end{solution}

\begin{problem}
已知两点 \(M\left(1,\frac{5}{4}\right)\)，\(N\left(-4,-\frac{5}{4}\right)\)，给出下列曲线方程：

\begin{circlelist}
\item \(4x+2y-1=0\)；
\item \(x^{2}+y^{2}=3\)；
\item \(\frac{x^{2}}{2}+y^{2}=1\)；
\item \(\frac{x^{2}}{2}-y^{2}=1\).
\end{circlelist}

在曲线上存在点 \(P\) 满足 \(\left|MP\right|=\left|NP\right|\) 的所有曲线方程是（\quad）

\choices
    {\(\circled{1}\circled{3}\)}
    {\(\circled{2}\circled{4}\)}
    {\(\circled{1}\circled{2}\circled{3}\)}
    {\(\circled{2}\circled{3}\circled{4}\)}
\end{problem}

\begin{answer}
D
\end{answer}

\begin{solution}
点 \(P(x,y)\) 满足 \(\left|MP\right|=\left|NP\right|\) 的轨迹是线段 \(MN\) 的垂直平分线. 由
\[
(x-1)^2+\left(y-\frac{5}{4}\right)^2=(x+4)^2+\left(y+\frac{5}{4}\right)^2
\]
得垂直平分线方程为 \(2x+y+3=0\)，即 \(y=-2x-3\).

曲线 \(\circled{1}\) 的方程 \(4x+2y-1=0\) 与该直线平行且不重合，没有交点. 对曲线 \(\circled{2}\)，代入得
\(x^2+(2x+3)^2=3\)，即 \(5x^2+12x+6=0\)，判别式为 \(24>0\)，有交点. 对曲线 \(\circled{3}\)，代入得
\(\frac{x^2}{2}+(2x+3)^2=1\)，即 \((3x+4)^2=0\)，有交点. 对曲线 \(\circled{4}\)，代入得
\(\frac{x^2}{2}-(2x+3)^2=1\)，即 \(7x^2+24x+20=0\)，判别式为 \(16>0\)，有交点. 因而所有符合条件的曲线为 \(\circled{2}\)、\(\circled{3}\)、\(\circled{4}\)，故选 D.
\end{solution}

\begin{problem}
某电脑用户计划使用不超过 \(500\) 元的资金购买单价分别为 \(60\) 元，\(70\) 元的单片软件和盒装磁盘，根据需要，软件至少买 \(3\) 片，磁盘至少买 \(2\) 盒，则不同的选购方式共有（\quad）

\choices
    {\(5\text{ 种}\)}
    {\(6\text{ 种}\)}
    {\(7\text{ 种}\)}
    {\(8\text{ 种}\)}
\end{problem}

\begin{answer}
C
\end{answer}

\begin{solution}
设软件购买 \(x\) 片，磁盘购买 \(y\) 盒，则 \(x\)，\(y\) 为正整数且 \(x\geqslant3\)，\(y\geqslant2\)，并满足
\(60x+70y\leqslant500\).
令 \(x'=x-3\)，\(y'=y-2\)，则 \(x'\)，\(y'\) 为非负整数，且
\(6x'+7y'\leqslant18\).
当 \(y'=0\) 时，\(x'=0\)，\(1\)，\(2\)，\(3\)，有 \(4\) 种；当 \(y'=1\) 时，\(x'=0\)，\(1\)，有 \(2\) 种；当 \(y'=2\) 时，只有 \(x'=0\)，有 \(1\) 种；\(y'\geqslant3\) 时不可能. 所以共有 \(4+2+1=7\) 种，故选 C.
\end{solution}

\section{填空题}

\begin{problem}
设椭圆 \(\frac{x^2}{a^2}+\frac{y^2}{b^2}=1\)（\(a>b>0\)）的右焦点为 \(F_1\)，右准线为 \(l_1\)，若过 \(F_1\) 且垂直于 \(x\) 轴的弦长等于点 \(F_1\) 到 \(l_1\) 的距离，则椭圆的离心率是\fillinblank{}.
\end{problem}

\begin{answer}
\(\frac{1}{2}\)
\end{answer}

\begin{solution}
设椭圆焦距的一半为 \(c\)，则 \(c^2=a^2-b^2\)，右焦点为 \(F_1(c,0)\)，右准线方程为 \(x=\frac{a^2}{c}\).
过 \(F_1\) 且垂直于 \(x\) 轴的弦所在直线为 \(x=c\). 代入椭圆方程得弦长
\(2\frac{b^2}{a}\). 点 \(F_1\) 到右准线的距离为
\(\frac{a^2}{c}-c=\frac{b^2}{c}\).
由题意 \(2\frac{b^2}{a}=\frac{b^2}{c}\)，所以 \(c=\frac{a}{2}\)，椭圆的离心率为 \(e=\frac{c}{a}=\frac{1}{2}\).
\end{solution}

\begin{problem}
在一块并排 \(10\text{ 垄}\) 的田地中，选择 \(2\text{ 垄}\) 分别种植 A，B 两种作物，每种作物种植一垄，为有利于作物生长，要求 A，B 两种作物的间隔不小于 \(6\text{ 垄}\)，则不同的选垄方法共有\fillinblank{}\(\text{种}\).
\end{problem}

\begin{answer}
12
\end{answer}

\begin{solution}
把 \(10\) 垄从左到右编号为 \(1\)，\(2\)，\(\ldots\)，\(10\). 若两种作物分别种在第 \(i\)，\(j\) 垄，则它们之间至少有 \(6\) 垄，故 \(|i-j|\geqslant7\).
无序地选择两垄时，间距为 \(7\)，\(8\)，\(9\) 的选法数分别为 \(3\)，\(2\)，\(1\)，共有 \(6\) 对. 由于 \(A\)，\(B\) 两种作物不同，每一对垄有 \(2\) 种分配方式，所以共有 \(2(3+2+1)=12\) 种.
\end{solution}

\begin{problem}
若正数 \(a\)，\(b\) 满足 \(ab = a + b + 3\)，则 \(ab\) 的取值范围是\fillinblank{}.
\end{problem}

\begin{answer}
\([9,+\infty)\)
\end{answer}

\begin{solution}
设 \(p=ab\). 由 \(ab=a+b+3\) 得 \(a+b=p-3\). 因为 \(a\)，\(b\) 是方程
\(u^2-(p-3)u+p=0\) 的两个正根，所以其判别式必须非负：
\((p-3)^2-4p=(p-1)(p-9)\geqslant0\).
又 \(p=ab=a+b+3>3\)，故只能有 \(p\geqslant9\).
当 \(p\geqslant9\) 时，方程的两个根为正数，且它们满足原条件，所以取值范围为 \([9,+\infty)\).
\end{solution}

\begin{problem}
\(\alpha\)，\(\beta\) 是两个不同的平面，\(m\)，\(n\) 是平面 \(\alpha\) 及 \(\beta\) 之外的两条不同直线. 给出四个论断：

\begin{circlelist}
\item \(m\perp n\)；
\item \(\alpha\perp\beta\)；
\item \(n\perp\beta\)；
\item \(m\perp\alpha\).
\end{circlelist}

以其中三个论断作为条件，余下一个论断作为结论，写出你认为正确的一个命题：\fillinblank{}.
\end{problem}

\begin{answer}
\(m\perp\alpha\)，\(n\perp\beta\)，\(\alpha\perp\beta\) \(\Longrightarrow\) \(m\perp n\)
\end{answer}

\begin{solution}
取平面 \(\alpha\)，\(\beta\) 的法向量分别为 \(\bs{p}\)，\(\bs{q}\)，并取直线 \(m\)，\(n\) 的方向向量分别为 \(\bs{u}\)，\(\bs{v}\). 由 \(m\perp\alpha\)，直线 \(m\) 的方向与平面 \(\alpha\) 的法向量平行，故 \(\bs{u}=\lambda\bs{p}\)；同理，由 \(n\perp\beta\)，有 \(\bs{v}=\mu\bs{q}\)，其中 \(\lambda\)，\(\mu\) 均不为零. 由 \(\alpha\perp\beta\)，两平面的法向量垂直，即 \(\bs{p}\cdot\bs{q}=0\). 因而
\[
\bs{u}\cdot\bs{v}=\lambda\mu\left(\bs{p}\cdot\bs{q}\right)=0
\]
所以直线 \(m\)，\(n\) 的方向互相垂直，即 \(m\perp n\). 因此，取第 2，3，4 个论断为条件，第 1 个论断为结论，得到的命题正确.
\end{solution}

\section{解答题}

\begin{problem}
解不等式：\(\sqrt{3\log_{a}x-2}<2\log_{a}x-1\)（\(a>0\)，\(a\neq 1\)）.
\end{problem}

\begin{solution}
令 \(t=\log_{a}x\). 根式有意义要求 \(3t-2\geqslant 0\). 又因为左边非负，而不等式为严格小于关系，所以右边必须为正，即 \(2t-1>0\). 因而定义域条件可合并为 \(t\geqslant\frac{2}{3}\). 在此条件下两边均非负，可以平方，原不等式等价于 \(3t-2<(2t-1)^2\)，即 \(4t^2-7t+3>0\)，也就是 \((4t-3)(t-1)>0\). 所以 \(t\in\left[\frac{2}{3},\frac{3}{4}\right)\cup(1,+\infty)\).

当 \(a>1\) 时，\(\log_a x\) 是增函数，故 \(x\in\left[a^{\frac{2}{3}},a^{\frac{3}{4}}\right)\cup(a,+\infty)\). 当 \(0<a<1\) 时，\(\log_a x\) 是减函数，故 \(x\in(0,a)\cup\left(a^{\frac{3}{4}},a^{\frac{2}{3}}\right]\).
\end{solution}

\begin{problem}
设复数 \(z=3\cos\theta+\i\cdot 2\sin\theta\)，求函数 \(y=\theta-\arg z\)（\(0<\theta<\frac{\pi}{2}\)）的最大值以及对应的 \(\theta\) 值.
\end{problem}

\begin{solution}
因为 \(0<\theta<\frac{\pi}{2}\)，复数 \(z\) 的实部和虚部均为正，所以
\(0<\arg z<\frac{\pi}{2}\)，且
\[
\tan(\arg z)=\frac{2\sin\theta}{3\cos\theta}=\frac{2}{3}\tan\theta
\]
令 \(t=\tan\theta>0\). 由 \(y=\theta-\arg z\) 以及正切差角公式，得
\[
\tan y=\frac{t-\frac{2}{3}t}{1+\frac{2}{3}t^2}=\frac{1}{3/t+2t}
\]
又 \(0<\arg z<\theta<\frac{\pi}{2}\)，所以 \(0<y<\frac{\pi}{2}\)，正切函数在此区间上严格递增. 由均值不等式，
\[
\frac{3}{t}+2t\geqslant2\sqrt{6}
\]
等号当且仅当 \(\frac{3}{t}=2t\)，即 \(t=\sqrt{\frac{3}{2}}=\frac{\sqrt{6}}{2}\). 因此
\[
\tan y\leqslant\frac{1}{2\sqrt{6}}=\frac{\sqrt{6}}{12}
\]
所以 \(y\) 的最大值为 \(\arctan\frac{\sqrt{6}}{12}\)，且对应
\(\theta=\arctan\frac{\sqrt{6}}{2}\).
\end{solution}

\begin{problem}
如图，已知正四棱柱 \(ABCD-A_{1}B_{1}C_{1}D_{1}\)，点 \(E\) 在棱 \(D_{1}D\) 上，截面 \(EAC\parallel D_{1}B\)，且面 \(EAC\) 与底面 \(ABCD\) 所成的角为 \(45^{\circ}\)，\(AB=a\).

\begin{enumerate}
\item 求截面 \(EAC\) 的面积；
\item 求异面直线 \(A_{1}B_{1}\) 与 \(AC\) 之间的距离；
\item 求三棱锥 \(B_{1}-EAC\) 的体积.
\end{enumerate}

\bitmapfigure[max width=0.62\linewidth]{img/1999/national_science/q21_fig1.png}
\end{problem}

\begin{solution}
\begin{enumerate}
\item 设 \(O=AC\cap BD\)，取 \(O\) 为原点，\(AC\) 为 \(x\) 轴，\(BD\) 为 \(y\) 轴，竖直向上为 \(z\) 轴，并令 \(s=\frac{a}{\sqrt{2}}\). 设正四棱柱的高为 \(h\)，则 \(A=(-s,0,0)\)，\(C=(s,0,0)\)，\(D=(0,-s,0)\)，\(B=(0,s,0)\)，\(D_1=(0,-s,h)\)，\(E=(0,-s,t)\)，其中 \(0\leqslant t\leqslant h\).

平面 \(EAC\) 的方程可写为 \(ty+sz=0\)，而 \(\overrightarrow{D_1B}=(0,2s,-h)\). 由 \(EAC\parallel D_1B\)，得 \((0,2s,-h)\cdot(0,t,s)=0\)，即 \(2st-sh=0\)，所以 \(h=2t\). 过 \(O\) 作垂直于 \(AC\) 的截面，该截面中平面 \(EAC\) 与底面的交线分别为 \(OE\) 和 \(OD\)，二面角为 \(45^{\circ}\)，故 \(\tan45^{\circ}=t/s\)，从而 \(t=s\)，于是 \(h=2s=\sqrt{2}a\).

此时 \(AC=2s=\sqrt{2}a\)，且点 \(E\) 到直线 \(AC\) 的距离为 \(\sqrt{s^2+t^2}=a\). 因而 \(S_{\triangle EAC}=\frac{1}{2}\cdot\sqrt{2}a\cdot a=\frac{\sqrt{2}}{2}a^2\).
\item 线段 \(A_1A\) 的两个端点分别在直线 \(A_1B_1\) 和直线 \(AC\) 上，且 \(A_1A\) 同时垂直于这两条直线，所以它是两异面直线的公垂线. 因而异面直线 \(A_1B_1\) 与 \(AC\) 之间的距离为 \(A_1A=h=\sqrt{2}a\).
\item 当 \(t=s\) 时，平面 \(EAC\) 的方程为 \(y+z=0\). 点 \(B_1=(0,s,h)=(0,s,2s)\) 到该平面的距离为 \(\frac{|s+2s|}{\sqrt{2}}=\frac{3a}{2}\). 因此三棱锥的高为 \(\frac{3a}{2}\)，从而 \(V_{B_1-EAC}=\frac{1}{3}\cdot\frac{\sqrt{2}}{2}a^2\cdot\frac{3a}{2}=\frac{\sqrt{2}}{4}a^3\).
\end{enumerate}
\end{solution}

\begin{problem}
如图为一台冷轧机的示意图. 冷轧机由若干对轧辊组成. 带钢从一端输入，经过各对轧辊逐步减薄后输出.

\begin{enumerate}
\item 输入带钢的厚度为 \(\alpha\)，输出带钢的厚度为 \(\beta\)，若每对轧辊的减薄率不超过 \(r_0\)，问冷轧机至少需要安装多少对轧辊（单对轧辊减薄率 \(=\frac{\text{输入该对的带钢厚度}-\text{从该对输出的带钢厚度}}{\text{输入该对的带钢厚度}}\)）？
\item 已知一台冷轧机共有 \(4\) 对减薄率为 \(20\%\) 的轧辊，所有轧辊周长均为 \(1600 \mathrm{~mm}\). 若第 \(k\) 对轧辊有缺陷，每滚动一周在带钢上压出一个疵点，在冷轧机输出的带钢上，疵点的间距为 \(L_k\). 为了便于检修，请计算 \(L_1\)，\(L_2\)，\(L_3\) 并填入下表（轧钢过程中，带钢宽度不变，且不考虑损耗）.
\end{enumerate}

\begin{center}
\begin{tabular}{|c|c|c|c|c|}
\hline
轧辊序号 & 1 & 2 & 3 & 4 \\
\hline
疵点间距 \(L_k/\mathrm{mm}\) &  &  &  & 1600 \\
\hline
\end{tabular}
\end{center}

\FigureLayoutDeclare{img/1999/national_science/q22_fig1.png}{16.0cm}{0bp 40bp 107bp 0bp}
\bitmapfigure[max width=0.9\linewidth]{img/1999/national_science/q22_fig1.png}
\end{problem}

\begin{solution}
\begin{enumerate}
\item 设安装 \(n\) 对轧辊. 在每对轧辊均达到允许的最大减薄率 \(r_0\) 时，厚度依次乘以 \(1-r_0\)，所以出口厚度的最小可能值为 \(\alpha(1-r_0)^n\). 在通常的减薄情形 \(0<\beta<\alpha\)、\(0<r_0<1\) 下，要使出口厚度不超过 \(\beta\)，必须且只需
\[
\alpha(1-r_0)^n\leqslant\beta
\]
由于 \(0<1-r_0<1\)，\(\lg(1-r_0)<0\)，取对数并反转不等号，得到
\[
n\geqslant\frac{\lg\beta-\lg\alpha}{\lg(1-r_0)}
\]
因此最少需要安装 \(\left\lceil\frac{\lg\beta-\lg\alpha}{\lg(1-r_0)}\right\rceil\) 对轧辊. 若 \(\beta\geqslant\alpha\)，则不安装轧辊已满足厚度要求；若 \(r_0=0\) 且 \(\beta<\alpha\)，则无法达到目标厚度.
\item 设带钢宽度为 \(w\)，输入厚度为 \(t_0\)，第 \(j\) 对轧辊出口处的厚度为 \(t_j\). 由于每对轧辊的减薄率为 \(20\%\)，有 \(t_j=0.8^j t_0\). 第 \(k\) 对轧辊出口处相邻疵点的间距是轧辊周长 \(1600\mathrm{~mm}\)，所以两疵点之间带钢的体积为 \(1600t_kw\). 这段带钢经过第 \(k+1\) 至第 \(4\) 对轧辊后，输出间距为 \(L_k\)，体积变为 \(L_kt_4w\). 由带钢宽度不变且轧制过程中不考虑损耗，得 \(1600t_kw=L_kt_4w\)，即 \(L_k=1600\frac{t_k}{t_4}=\frac{1600}{0.8^{4-k}}\). 因而 \(L_1=3125\mathrm{~mm}\)，\(L_2=2500\mathrm{~mm}\)，\(L_3=2000\mathrm{~mm}\)，且 \(L_4=1600\mathrm{~mm}\)，与表中已给数值一致.
\end{enumerate}
\end{solution}

\begin{problem}
已知函数 \(y=f(x)\) 的图象是自原点出发的一条折线. 当 \(n\leqslant y\leqslant n+1\)（\(n=0\)，\(1\)，\(2\)，\(\cdots\)）时，该图象是斜率为 \(b^n\) 的线段（其中正常数 \(b\neq 1\)），设数列 \(\{x_n\}\) 由 \(f(x_n)=n\)（\(n=1\)，\(2\)，\(\cdots\)）定义.

\begin{enumerate}
\item 求 \(x_1\)，\(x_2\) 和 \(x_n\) 的表达式；
\item 求 \(f(x)\) 的表达式，并写出其定义域；
\item 证明：\(y=f(x)\) 的图象与 \(y=x\) 的图象没有横坐标大于 \(1\) 的交点.
\end{enumerate}
\end{problem}

\begin{solution}
\begin{enumerate}
\item 当 \(0\leqslant y\leqslant1\) 时，图象从原点出发且斜率为 \(b^0=1\)，所以 \(x_1=1\). 当 \(1\leqslant y\leqslant2\) 时斜率为 \(b\)，故 \(x_2-x_1=1/b\)，从而 \(x_2=1+1/b\). 一般地，第 \(n\) 段线段的斜率为 \(b^n\)，所以 \(x_{n+1}-x_n=1/b^n\). 因此
\(x_n=\sum_{k=0}^{n-1}\frac{1}{b^k}=\frac{1-b^{-n}}{1-b^{-1}}=\frac{b-(1/b)^{n-1}}{b-1}\)，其中 \(n=1\)，\(2\)，\(\ldots\).
\item 当 \(0\leqslant x\leqslant1\) 时，\(f(x)=x\). 当 \(n\leqslant y\leqslant n+1\) 且 \(n\geqslant1\) 时，线段经过 \((x_n,n)\)，斜率为 \(b^n\)，故 \(f(x)=n+b^n(x-x_n)\)，其中 \(x_n\leqslant x\leqslant x_{n+1}\). 由 \(x_n=\sum_{k=0}^{n-1}b^{-k}\) 可知：当 \(b>1\) 时，\(x_n\) 递增并趋近于 \(\frac{b}{b-1}\)，所以定义域为 \(\left[0,\frac{b}{b-1}\right)\)；当 \(0<b<1\) 时，\(x_n\to+\infty\)，所以定义域为 \([0,+\infty)\).
\item 对于 \(b>1\)，在第一段 \(1<x\leqslant x_2\) 上，\(f(x)-x=(b-1)(x-1)>0\). 对于任意 \(n\geqslant2\)，因为 \(b^{-k}<1\)（\(k\geqslant1\)），所以 \(x_n<n\)，\(x_{n+1}<n+1\). 在线段 \([x_n,x_{n+1}]\) 上，\(f(x)-x\) 是关于 \(x\) 的一次函数，并且在两个端点都为正，故整个线段上均为正.

对于 \(0<b<1\)，在第一段 \(1<x\leqslant x_2\) 上，\(f(x)-x=(b-1)(x-1)<0\). 对于 \(n\geqslant2\)，因为 \(b^{-k}>1\)（\(k\geqslant1\)），所以 \(x_n>n\)，从而在 \([x_n,x_{n+1}]\) 的两个端点处均有 \(f(x)-x<0\). 由于差值在线段上是一次函数，故整个线段上均为负. 因而当 \(x>1\) 时总有 \(f(x)\neq x\)，即两图象没有横坐标大于 \(1\) 的交点.
\end{enumerate}
\end{solution}

\begin{problem}
如图，给出定点 \(A(a,0)\)（\(a>0\)）和直线 \(l\colon x=-1\). \(B\) 是直线 \(l\) 上的动点，\(\angle BOA\) 的角平分线交 \(AB\) 于点 \(C\). 求点 \(C\) 的轨迹方程，并讨论方程表示的曲线类型与 \(a\) 值的关系.

\bitmapfigure[max width=0.55\linewidth]{img/1999/national_science/q24_fig1.png}
\end{problem}

\begin{solution}
设动点 \(B=(-1,t)\)，记 \(r=OB=\sqrt{1+t^2}\). 在三角形 \(AOB\) 中，\(OC\) 是顶角的内角平分线，由角平分线定理得 \(AC:CB=AO:OB=a:r\). 因而 \(C=\frac{rA+aB}{a+r}\)，即 \(x=\frac{a(r-1)}{a+r}\)，\(y=\frac{at}{a+r}\).

由上式可得 \(a-x=\frac{a(a+1)}{a+r}\)，从而 \(r=\frac{a(x+1)}{a-x}\)，同时 \(t=\frac{(a+1)y}{a-x}\). 代入 \(r^2=1+t^2\)，得
\[
a^2(x+1)^2=(a-x)^2+(a+1)^2y^2
\]
整理得点 \(C\) 的轨迹方程
\[
(a-1)x^2+2ax-(a+1)y^2=0
\]
由参数表示还可知实际轨迹满足 \(0\leqslant x<a\). 反过来，先取二次曲线上满足 \(0<x<a\) 的一点 \((x,y)\)，分别令
\(r=\frac{a(x+1)}{a-x}\)，\(t=\frac{(a+1)y}{a-x}\).
此时 \(r>1\)，且由轨迹方程可得 \(r^2=1+t^2\). 取 \(B=(-1,t)\)，则 \(OB=r\)，三角形 \(AOB\) 非退化，代入角平分线定理的内分公式
\[
C=\frac{rA+aB}{a+r}
\]
即可得到原点为顶点的角平分线与 \(AB\) 的交点正是 \((x,y)\). 当 \(x=0\) 时，轨迹方程给出 \(y=0\)，对应 \(B=(-1,0)\) 时的退化交点 \(C=O\). 所以若把该退化情形计入，实际轨迹是该二次曲线中 \(0\leqslant x<a\) 的部分；若排除 \(B=(-1,0)\)，则相应取 \(0<x<a\).

当 \(a>1\) 时，方程可化为 \((a-1)\left(x+\frac{a}{a-1}\right)^2-(a+1)y^2=\frac{a^2}{a-1}\)，表示双曲线. 当 \(a=1\) 时，方程为 \(x=y^2\)，表示抛物线. 当 \(0<a<1\) 时，方程可化为 \((1-a)\left(x-\frac{a}{1-a}\right)^2+(a+1)y^2=\frac{a^2}{1-a}\)，表示椭圆.
\end{solution}
